% Section VI -- Conclusion. No number here that is not in Section IV.

\section{Conclusion}

We measured what differential privacy costs in energy when added to federated
learning, using hardware counters around each round on two datasets, with
UCI-HAR partitioned so each client holds one study participant's data.

The cost divides in two. A fixed per-round overhead --- $+45\%$ on HAR,
$+60\%$ on CIFAR-10, flat in $\varepsilon$ within our resolution --- comes
from per-sample gradients and clipping. Everything else acts through the
number of rounds required, and it is the larger term: reaching $0.70$ accuracy
on HAR costs $5.8\times$ the non-private energy at $\varepsilon = 4$ and is
unreachable at $\varepsilon \le 2$ regardless of expenditure. Separating them
matters more than quantifying either, since they respond to different
interventions.

Three rules follow. Derive the round budget from the privacy budget rather
than the task, since rounds consume budget and training to convergence is
self-defeating --- at $\varepsilon = 0.5$ on HAR, $83\%$ of a run's energy
falls after the accuracy peak. Choose the privacy budget and the accuracy
requirement together, because a budget implies a ceiling no computation
reaches. And when the private model is noise-limited, hyperparameter search is
not the lever.

Paid in joules, the cost inherits the carbon intensity of wherever training
runs: an $18.8\times$ range between Algeria's grid and Norway's for an
identical guarantee. A privacy--utility evaluation cannot represent this, and
it is awkward for federated learning, where the computation sits on whatever
grids the participating devices do.

We also report a measurement hazard. Instrumentation attached to a privacy
mechanism scales with the privacy parameter, so its cost mimics the effect
under study; deriving noise multipliers before measurement, verifying step
counts per round, and randomising run order are cheap controls we recommend
for any such study.

Two limits bound the claims: clients are virtual and share one GPU, so the
ratios transfer but the absolute joules do not, and our two datasets bound the
effects rather than establishing them. Characterising where the boundary
between those regimes falls --- whether it is set by task difficulty, by the
magnitude of the noise, or by their ratio --- is the question this work leaves
open.
